% ============================================================
% Chapter 6: Experimental Design
%   Section: Experiments on the Stages Before the Verdict
% Included from chapters/06-experimental-design/experimental-design.tex
% Source: docs/experiments.md; backend/src/evaluation/experiments.py;
% results in backend/data/evaluation/experiments/.
% ============================================================

\section{Experiments on the Stages Before the Verdict}
\label{sec:experiments-pre-verdict}

The harness measures what the fact-checker concludes. A second tool
measures what decides which articles reach it at all --- discovery, the
selection step, admission and the topic classifier --- which is where
the answer to RQ1 and part of RQ3 lie. One command runs each experiment:

\begin{lstlisting}
uv run python -m src.evaluation.experiments <experiment> [options]
\end{lstlisting}

and writes a JSON result and a readable report to a timestamped folder
under \path{backend/data/evaluation/experiments/}. Intervals are
percentile 95\% bootstrap intervals with a fixed seed and 2,000
resamples; an AUC is the probability that a random positive outranks a
random negative, 0.5 being a coin.

\begin{table}[ht]
\centering
\small
\caption{The experiments on the stages before the verdict.}
\label{tab:experiments}
\begin{tabularx}{\linewidth}{@{}lXX@{}}
\toprule
\textbf{Experiment} & \textbf{Question} & \textbf{Measures} \\
\midrule
\texttt{discovery} & How much does scoping to topic groups save, and
what does it still find? & Sources read, attempts, links, candidates by
language, seconds, each against the every-topic baseline \\
\texttt{selection} & How far does the editor agree with the AI
selection? & Precision, recall and Jaccard of the AI's picks against
what was sent; AUC of its score; assisted and blind rounds apart \\
\texttt{blind-ai} & The same, without the editor having seen the AI's
picks & Runs the AI on a round already sent \\
\texttt{admission} & What do the admission thresholds admit? &
Admission rate, and precision and recall against the editor's choice,
over a grid of topic and impact minimums; AUC of the impact score \\
\texttt{topics} & How often is the topic classifier right? & Top-1
topic, gold in the top three, top-1 group, share with no topic \\
\texttt{emphasis} & What does weighing the topic over positive impact
change? & Re-scores recorded rounds at each emphasis; their first ten
and twenty against labels \\
\bottomrule
\end{tabularx}
\end{table}

Two further studies were run once rather than as commands, to choose
values that are now constants of the code: the mission screen's margin
(527 live candidates read by hand) and the ranking's weights (117
candidates labelled before any score was computed, with every reading
kept so other weights can be tried without embedding anything again).
Their data are committed with the experiments.

\paragraph{Labels.} The candidate labels these experiments are scored
against --- clearly positive, worth offering, not worth offering, and
whether a page is an article at all --- were assigned from each
candidate's title, summary and link. The labels of the 6 and
9~October rounds were assigned by the AI labelling pass (Claude) and
have not yet been checked by a person; the P1 items of
Table~\ref{tab:priorities} require the author to label at least two
rounds and to run ten blind selection rounds before RQ1 and RQ3 are
answered. Feeds change through the day, so settings are compared
within one run and repeated on several days for an interval.
